\chapter{定積分の対称性}\label{章:定積分の対称性}

不定積分では,微分すると与えられた関数になる式を探した.定積分でも原始関数が見つかれば計算できるが,定積分にはもう1つ情報がある.どこからどこまで積分するかという区間である.たとえば,同じ関数でも区間が原点の左右に同じ長さだけ延びていれば,左右の値を組にして考えられる.区間の端が互いに逆数なら,逆数の置換によって両端を交換できる.

この章では,区間の両端を交換する置換をして,もとと同じ値の積分を作る.置換後の積分はすぐには計算せず,置換前の積分と並べてみる.2つを加えたとき,分母が消えることもあれば,変数を掛けた部分が定数に変わることもある.そうなれば,もとの積分を未知数とする方程式を解くだけで定積分の値が決まる.原始関数を求めるには手間がかかる式でも,短い計算で済むのである.

ここで1つ断りをいれておこう.以下では,特に断らない限り,被積分関数は積分する閉区間で連続であるとする.途中に分母が $0$ になる点がある式に,対称性だけを見て同じ議論を当てはめることはできない.そのような点を含む積分は,\chapterref{章:近似と広義積分}で定義を確かめてから扱う.

\section{偶関数と奇関数}\label{節:奇偶の対称性}

\begin{definition}{偶関数と奇関数}
原点に関して対称な範囲で定義された関数 $f$ について,すべての $x$ で
\[
f(-x)=f(x)
\]
が成り立つとき,\,$f$ を偶関数という.また,すべての $x$ で
\[
f(-x)=-f(x)
\]
が成り立つとき,\,$f$ を奇関数という.
\end{definition}

偶関数のグラフは $y$ 軸について対称であり,奇関数のグラフは原点について対称である.\,$x^2$ と $\cos x$ は偶関数,\,$x^3$ と $\sin x$ は奇関数である.ただし,式の中に $\sin$ があれば奇関数になるわけではない.\,$\sin(x^2)$ では $x$ を $-x$ に変えても $x^2$ が変わらないので偶関数である.式全体に $-x$ を代入することが,確実な判定になる.

区間が原点について対称なとき,この対称性を積分に使えないだろうか.$a>0$ とし,\,$[-a,a]$ での積分を,区間 $[0,a]$ での積分に帰着させることを目標とする.

$[-a,a]$ の積分を $0$ で分ける.左半分で $x=-t$ と置くと,\,$x=-a,0$ に対して $t=a,0$ であり,\,$\dd x=-\dd t$ だから,
\[
\begin{aligned}
\int_{-a}^{0}f(x)\dd x
 &=\int_a^0 f(-t)(-\dd t)\\
 &=\int_0^a f(-t)\dd t
\end{aligned}
\]
となる.積分変数の文字を $x$ に戻して右半分と加えれば,
\begin{equation}\label{式:左右の積分}
\int_{-a}^a f(x)\dd x
 =\int_0^a\{f(x)+f(-x)\}\dd x
\end{equation}
を得る.これは偶関数でも奇関数でもない場合にも使える式である.偶関数なら中括弧が $2f(x)$ になり,奇関数なら $0$ になる.

\begin{theorem}{対称区間での偶関数・奇関数の積分}
$a>0$ とする.\,$f$ が偶関数なら
\[
\int_{-a}^a f(x)\dd x=2\int_0^a f(x)\dd x,
\]
$f$ が奇関数なら
\[
\int_{-a}^a f(x)\dd x=0
\]
である.
\end{theorem}

図\ref{図:奇偶の積分}の右図では,左右の部分の幾何学的な面積は等しい.しかし左側の定積分は負,右側は正である.奇関数の積分が $0$ という結論は,図形の面積が $0$ という意味ではない.面積を求めるなら絶対値を積分する必要がある.

\begin{figure}[htbp]
\centering
\begin{tikzpicture}[scale=1.05,>=Stealth]
\begin{scope}[xshift=-3.1cm]
\fill[blue!12] plot[domain=-1.35:1.35,samples=50](\x,{0.65*\x*\x})--(1.35,0)--(-1.35,0)--cycle;
\draw[->](-1.7,0)--(1.8,0)node[right]{$x$};
\draw[->](0,-0.2)--(0,1.7)node[above]{$y$};
\draw[blue!70!black,thick] plot[domain=-1.45:1.45,samples=60](\x,{0.65*\x*\x});
\draw[dashed](-1.35,0)--(-1.35,1.185);
\draw[dashed](1.35,0)--(1.35,1.185);
\node[below]at(-1.35,0){$-a$};\node[below]at(1.35,0){$a$};
\node at(0,-0.65){偶関数：左右を加える};
\end{scope}
\begin{scope}[xshift=1.8cm]
\fill[blue!12](0,0)--plot[domain=0:1.35,samples=40](\x,{0.55*\x*\x*\x})--(1.35,0)--cycle;
\fill[orange!18](0,0)--plot[domain=0:-1.35,samples=40](\x,{0.55*\x*\x*\x})--(-1.35,0)--cycle;
\draw[->](-1.7,0)--(1.8,0)node[right]{$x$};
\draw[->](0,-1.55)--(0,1.7)node[above]{$y$};
\draw[blue!70!black,thick] plot[domain=-1.4:1.4,samples=60](\x,{0.55*\x*\x*\x});
\draw[dashed](-1.35,0)--(-1.35,-1.353);
\draw[dashed](1.35,0)--(1.35,1.353);
\node[above]at(-1.35,0){$-a$};\node[below]at(1.35,0){$a$};
\node at(0,-1.95){奇関数：符号付きの量が打ち消す};
\end{scope}
\end{tikzpicture}
\caption{対称性と,定積分の符号の関係.曲線はそれぞれ $x^2$ 型,\,$x^3$ 型である.}\label{図:奇偶の積分}
\end{figure}

\begin{bookpoint}
\textbf{POINT}\par
$x$ を $-x$ に変えて各項の偶奇を調べ,奇関数の項は $0$ とし,偶関数の項は $[0,a]$ での積分の2倍にする.
\end{bookpoint}

\begin{example}{偶奇による分解}
次の定積分を求めよ.
\[
I=\int_{-2}^{2}\left(x^3\e^{-x^2}+2x^2+1\right)\dd x.
\]
\end{example}

\paragraphbf{方針}
被積分関数を $x^3\e^{-x^2}$ と $2x^2+1$ に分け,それぞれの偶奇を調べる.

\paragraphbf{解答}
$I=\frac{44}{3}$

\paragraphbf{解法}
$x$ を $-x$ に変えると,\,$x^3$ の符号だけが変わり,\,$\e^{-x^2}$ は変わらない.したがって積 $x^3\e^{-x^2}$ は奇関数で,\,$[-2,2]$ での積分は $0$ である.残った $2x^2+1$ は偶関数だから,
\[
\begin{aligned}
I&=2\int_0^2(2x^2+1)\dd x\\
 &=2\left[\frac{2x^3}{3}+x\right]_0^2\\
 &=2\left(\frac{16}{3}+2\right)\\
 &=\frac{44}{3}
\end{aligned}
\]
である.\,$x^3\e^{-x^2}$ の原始関数は求めていない.また,式全体が偶関数か奇関数かを決める必要もなかった.和を分けて,対称性が分かる部分にだけ使えばよい.

\paragraphbf{偶関数部分と奇関数部分}
さらに,一般の $f$ に対して
\[
f_{\mathrm e}(x)=\frac{f(x)+f(-x)}{2},\qquad
f_{\mathrm o}(x)=\frac{f(x)-f(-x)}{2}
\]
と置くと,\,$f=f_{\mathrm e}+f_{\mathrm o}$ であり,前者は偶関数,後者は奇関数になる.よって対称区間の定積分に寄与するのは $f_{\mathrm e}$ の部分だけである.実際,式\eqref{式:左右の積分}の中括弧は $f(x)+f(-x)=2f_{\mathrm e}(x)$ であり,この式は $f$ から偶関数の部分を取り出して $2\int_0^a f_{\mathrm e}(x)\dd x$ を計算していたのである.
%
\section{区間の反転}\label{節:区間の反転}

前節では,区間が原点について対称であった.では,原点を中心としない区間 $[a,b]$ でも,左右の値を組にして考えられないだろうか.本節では特に,区間の中央について対称な関数 $g$,すなわち
\[
g(a+b-x)=g(x)
\]
を満たす $g$ に一次因子 $x$ が掛かった積分
\[
\int_a^b xg(x)\dd x
\]
を考える.この積分から $x$ を取り除ければ,あとは $\int_a^b g(x)\dd x$ を計算するだけで済む.

少々発想しにくい変形であるため,結果からお見せしよう.
\begin{equation}\label{式:対称な重みの一次因子}
\int_a^b xg(x)\dd x
 =\frac{a+b}{2}\int_a^b g(x)\dd x
\end{equation}
が成り立つ.つまり,積分の中の $x$ を区間の中央 $\frac{a+b}{2}$ に置き換えてよいのだ.

これを示すために,区間 $[a,b]$ の両端を交換する置換 $x=a+b-t$ を考える.対応は
\[
x=a\longleftrightarrow t=b,\qquad
x=b\longleftrightarrow t=a
\]
であり,\,$\dd x=-\dd t$ だから,一般の関数 $f$ について
\[
\begin{aligned}
\int_a^b f(x)\dd x
 &=\int_b^a f(a+b-t)(-\dd t)\\
 &=\int_a^b f(a+b-t)\dd t
\end{aligned}
\]
となる.置換の直後は積分の上下端が逆になるが,微分の負号と合わせて向きを戻している.積分変数の文字を $x$ に戻せば,
\begin{equation}\label{式:区間反転}
\int_a^b f(x)\dd x=\int_a^b f(a+b-x)\dd x
\end{equation}
を得る.

$I=\int_a^b xg(x)\dd x$ と置き,式\eqref{式:区間反転}を $f(x)=xg(x)$ に使う.\,$g(a+b-x)=g(x)$ だから
\[
I=\int_a^b(a+b-x)g(x)\dd x
\]
である.もとの式と加えると,\,$x+(a+b-x)=a+b$ なので
\[
\begin{aligned}
2I&=\int_a^b\{x+(a+b-x)\}g(x)\dd x\\
 &=(a+b)\int_a^b g(x)\dd x
\end{aligned}
\]
となる.両辺を $2$ で割れば式\eqref{式:対称な重みの一次因子}を得る.

ここで行ったのは,反転した積分ともとの積分を加えて $2$ で割ることである.一般の $f$ について書けば
\begin{equation}\label{式:反転平均}
\int_a^b f(x)\dd x
 =\frac12\int_a^b\{f(x)+f(a+b-x)\}\dd x
\end{equation}
となる.式\eqref{式:反転平均}の中括弧が定数 $k$ になれば,答えは $\frac12k(b-a)$ である.もとの式を個別に積分できるかどうかは,この計算では問題にならない.次節以降では,中括弧が定数になる例を扱う.

\begin{bookpoint}
\textbf{POINT}\par
$x=a+b-t$ と置いて反転した積分を作り,もとの積分と加える.\,$g(a+b-x)=g(x)$ なら,\,$x$ を区間の中央 $\frac{a+b}{2}$ に置き換えられる.
\end{bookpoint}

\begin{example}{反転による計算}
次の定積分を求めよ.
\[
I=\int_0^\pi x\sin^3x\dd x.
\]
\end{example}

\paragraphbf{方針}
$\sin(\pi-x)=\sin x$ より,\,$\sin^3x$ は区間の中央 $\frac{\pi}{2}$ について対称である.反転で $x$ を $\frac{\pi}{2}$ に置き換え,残りは $u=\cos x$ と置く.

\paragraphbf{解答}
$I=\frac{2\pi}{3}$

\paragraphbf{解法}
$x$ と三角関数の積を見ると部分積分を考えたくなるだろう.しかしこの区間では $\sin(\pi-x)=\sin x$ であり,\,$g(x)=\sin^3x$ は $g(\pi-x)=g(x)$ を満たす.式\eqref{式:対称な重みの一次因子}で $a=0$,\,$b=\pi$ とすれば
\[
I=\frac{\pi}{2}\int_0^\pi\sin^3x\dd x
\]
となる.残りは $\sin^3x=(1-\cos^2x)\sin x$ として $u=\cos x$ と置く.\,$x=0,\pi$ に対して $u=1,-1$ であり,\,$\sin x\dd x=-\dd u$ だから
\[
\begin{aligned}
\int_0^\pi\sin^3x\dd x
 &=\int_1^{-1}(1-u^2)(-\dd u)\\
 &=\int_{-1}^1(1-u^2)\dd u\\
 &=\left[u-\frac{u^3}{3}\right]_{-1}^1\\
 &=\frac43
\end{aligned}
\]
である.したがって
\[
I=\frac{\pi}{2}\cdot\frac43=\frac{2\pi}{3}
\]
を得る.反転で $x$ を取り除き,置換で三角関数を取り除くという2段階で計算できた.

\paragraphbf{反転の発想}
以下,\,$x$ を $a+b-x$ に移す変換を $x\mapsto a+b-x$ のように矢印で書く.反転 $x\mapsto a+b-x$ は,前節の $x\mapsto-x$ を区間の中央 $c=\frac{a+b}{2}$ のまわりで行ったものである.実際,\,$x=c+s$ と置けば $a+b-x=2c-(c+s)=c-s$ であり,反転は $s\mapsto-s$ にほかならない.

この見方で式\eqref{式:対称な重みの一次因子}をもう一度導いてみよう.\,$h=\frac{b-a}{2}$ とし,\,$x=c+s$ と置くと,区間は $-h\le s\le h$ となる.\,$G(s)=g(c+s)$ と置けば,\,$g(a+b-x)=g(x)$ は $G(-s)=G(s)$,すなわち $G$ が偶関数であることを表す.したがって
\[
\begin{aligned}
\int_a^b xg(x)\dd x
 &=\int_{-h}^h(c+s)G(s)\dd s\\
 &=c\int_{-h}^hG(s)\dd s+\int_{-h}^h sG(s)\dd s
\end{aligned}
\]
となる.\,$sG(s)$ は奇関数なので第2項は $0$ であり,第1項は $c\int_a^b g(x)\dd x$ に等しい.中央からの距離が等しい2点では $g$ の値が等しく,その2点の $x$ 座標の平均は常に区間の中央 $c$ になる.式\eqref{式:対称な重みの一次因子}は,このことを積分で表したものである.

\paragraphbf{反転が使えない場合}
式\eqref{式:対称な重みの一次因子}を使う際には,\,$xg(x)$ 全体が対称なのではなく,\,$g$ が対称であることを確認する.たとえば $\int_0^1 x\e^x\dd x$ では $\e^{1-x}=\e^x$ が成り立たないので,\,$x$ を区間の中央 $\frac12$ に置き換えることはできない.

また,\,$x^2g(x)$ まで同じように中央の2乗に置き換えられるわけでもない.\,$x^2+(a+b-x)^2$ は $x$ を含むからである.反転して足すことで何が消えるのかを,その都度式で確かめる必要がある.
%
\section{正弦と余弦}\label{節:正弦余弦の交換}

前節の反転を区間 $\left[0,\frac{\pi}{2}\right]$ で使うと何が起こるだろうか.この区間での反転は,\,$x$ を $\frac{\pi}{2}-x$ に変える操作である.このとき
\[
\sin\left(\frac{\pi}{2}-x\right)=\cos x,\qquad
\cos\left(\frac{\pi}{2}-x\right)=\sin x
\]
なので,正弦と余弦が入れ替わる.この入れ替えが役立つのは,\,$q$ を $[0,1]$ 上の連続関数として
\[
I=\int_0^{\frac{\pi}{2}}
\frac{q(\sin x)}{q(\sin x)+q(\cos x)}\dd x
\]
の形の積分である.ただし,分母 $q(\sin x)+q(\cos x)$ は区間内で $0$ にならないとする.反転した積分ともとの積分を加えて被積分関数を定数 $1$ にできれば,積分は区間の長さだけで決まる.

反転すると,分子の $q(\sin x)$ は $q(\cos x)$ に変わり,分母は2つの項が入れ替わるだけで変わらない.したがって
\[
I=\int_0^{\frac{\pi}{2}}
\frac{q(\cos x)}{q(\sin x)+q(\cos x)}\dd x
\]
である.もとの式と加えれば,分子の和が分母そのものになる.
\[
\begin{aligned}
2I&=\int_0^{\frac{\pi}{2}}
\frac{q(\sin x)+q(\cos x)}{q(\sin x)+q(\cos x)}\dd x\\
 &=\int_0^{\frac{\pi}{2}}1\dd x\\
 &=\frac{\pi}{2}
\end{aligned}
\]
よって
\[
\int_0^{\frac{\pi}{2}}
\frac{q(\sin x)}{q(\sin x)+q(\cos x)}\dd x
 =\frac{\pi}{4}
\]
を得る.\,$q$ の具体的な形は一度も使っていない.重要なのは,反転で分子の2つの候補が交換されることなのだ.

\begin{bookpoint}
\textbf{POINT}\par
$x=\frac{\pi}{2}-t$ と置いて正弦と余弦を入れ替え,もとの積分と加えて分子の和を分母に一致させる.
\end{bookpoint}

\begin{example}{正弦と余弦の交換}
次の定積分を求めよ.
\[
I=\int_0^{\frac{\pi}{2}}
\frac{\sin^4x}{\sin^4x+\cos^4x}\dd x.
\]
\end{example}

\paragraphbf{方針}
本節の式で $q(s)=s^4$ とした場合である.反転で分子の $\sin^4x$ が $\cos^4x$ に変わり,2つの分子の和がちょうど分母 $\sin^4x+\cos^4x$ になる.

\paragraphbf{解答}
$I=\frac{\pi}{4}$

\paragraphbf{解法}
$\sin x$ と $\cos x$ は同時に $0$ にならないので,分母は区間全体で正である.\,$x=\frac{\pi}{2}-t$ と置くと,分子は $\cos^4t$ に変わり,分母は変わらない.積分変数の文字を $x$ に戻せば
\[
I=\int_0^{\frac{\pi}{2}}
\frac{\cos^4x}{\sin^4x+\cos^4x}\dd x
\]
である.2つを加えれば,分子の和が分母そのものになる.
\[
\begin{aligned}
2I&=\int_0^{\frac{\pi}{2}}1\dd x\\
 &=\frac{\pi}{2}
\end{aligned}
\]
よって $I=\frac{\pi}{4}$ である.これは上の一般形で $q(t)=t^4$ とした場合である.\,$\sin^4x$ の次数を下げたり,分母を展開したりする必要はない.

なお,一般形は $q(t)=t^4$ に限らない.たとえば $q(t)=t^2+1$ でも,分母は $\sin^2x+\cos^2x+2=3$ で $0$ にならず,同じ値 $\frac{\pi}{4}$ になる.

\paragraphbf{図形的解釈}
図\ref{図:相補重み}は,例題の被積分関数を $w(x)$ としたグラフである.左右を反転したグラフ $w(\frac{\pi}{2}-x)$ と,同じ $x$ における高さを足すと $1$ になる.2つのグラフの下の面積は等しく,その和は縦 $1$,横 $\frac{\pi}{2}$ の長方形の面積である.したがって $w(x)$ のグラフは,この長方形の面積を2等分しているのである.

\begin{figure}[htbp]
\centering
\begin{tikzpicture}[x=3.5cm,y=2.3cm,>=Stealth]
\fill[blue!10](0,0)--plot[domain=0:1.570796,samples=100](\x,{pow(sin(\x r),4)/(pow(sin(\x r),4)+pow(cos(\x r),4))})--(1.570796,0)--cycle;
\draw[->](-0.08,0)--(1.73,0)node[right]{$x$};
\draw[->](0,-0.08)--(0,1.2)node[above]{$y$};
\draw[dashed](0,1)--(1.570796,1)--(1.570796,0);
\draw[dashed](0.785398,0)--(0.785398,1);
\draw[blue!75!black,thick]plot[domain=0:1.570796,samples=100](\x,{pow(sin(\x r),4)/(pow(sin(\x r),4)+pow(cos(\x r),4))});
\draw[orange!80!black,thick,dashed]plot[domain=0:1.570796,samples=100](\x,{pow(cos(\x r),4)/(pow(sin(\x r),4)+pow(cos(\x r),4))});
\node[left]at(0,1){$1$};\node[below]at(0,0){$0$};
\node[below]at(0.785398,0){$\frac{\pi}{4}$};
\node[below]at(1.570796,0){$\frac{\pi}{2}$};
\node[blue!75!black]at(1.15,0.65){$w(x)$};
\node[orange!80!black]at(0.48,1.26){$w(\frac{\pi}{2}-x)$};
\end{tikzpicture}
\caption{反転した2つの重みは,同じ場所で足すと $1$ になる.}\label{図:相補重み}
\end{figure}
%
\section{周期性}\label{節:周期関数の積分}

$T>0$ に対して $f(x+T)=f(x)$ がすべての $x$ で成り立つとき,\,$T$ を $f$ の周期という.周期という言葉だけから,\,$T$ は最小の正の周期だとは限らない.\,$\sin x$ の周期として $2\pi$ も $4\pi$ も使える.

周期関数をちょうど一周期分だけ積分した値は,始点によらない.まず $0\leq a\leq T$ の場合,
\[
\int_a^{a+T}f(x)\dd x
 =\int_a^T f(x)\dd x+\int_T^{a+T}f(x)\dd x.
\]
第二項で $x=t+T$ と置くと $f(t+T)=f(t)$ だから,
\[
\int_T^{a+T}f(x)\dd x=\int_0^a f(t)\dd t.
\]
二つを足すと
\begin{equation}\label{式:一周期の積分}
\int_a^{a+T}f(x)\dd x=\int_0^T f(x)\dd x
\end{equation}
になる.一般の $a$ も,整数個の $T$ だけ平行移動すれば始点を $[0,T)$ に移せるため,同じ結論になる.区間を長さ $T$ の部分に分ければ,正の整数 $n$ に対して
\[
\int_a^{a+nT}f(x)\dd x=n\int_0^T f(x)\dd x
\]
も得られる.

周期があるというだけでは,一周期分の積分が零とは限らない.たとえば $\sin^2x$ は周期 $\pi$ を持つが,一周期の積分は $\frac{\pi}{2}$ である.符号付きの部分が打ち消し合うか,積分区間が周期の整数倍になっているかを別々に確認する.

\begin{example}{絶対値がある周期関数}
次の定積分を求めよ.
\[
I=\int_{-\frac{\pi}{6}}^{\frac{13\pi}{6}}|\sin x|\dd x.
\]
\end{example}

$\sin(x+\pi)=-\sin x$ なので,絶対値を取ると周期は $\pi$ になる.一周期分は
\[
\int_0^\pi|\sin x|\dd x=\int_0^\pi\sin x\dd x=2.
\]
与えられた区間の長さは $\frac{7\pi}{3}$ であり,\,$2\pi$ と $\frac{\pi}{3}$ に分けられる.まず $-\frac{\pi}{6}$ から $\frac{11\pi}{6}$ までは二周期分なので積分は $4$ である.残りでは $x=t+2\pi$ と置き,
\[
\int_{\frac{11\pi}{6}}^{\frac{13\pi}{6}}|\sin x|\dd x
 =\int_{-\frac{\pi}{6}}^{\frac{\pi}{6}}|\sin t|\dd t.
\]
$|\sin t|$ は偶関数で,\,$0\leq t\leq\frac{\pi}{6}$ では $\sin t\geq0$ だから,
\[
I=4+2\int_0^{\frac{\pi}{6}}\sin t\dd t
 =4+2\left(1-\frac{\sqrt3}{2}\right)
 =6-\sqrt3.
\]
周期によって長い区間を短くし,偶奇によって符号の変わる区間を半分にした.周期の整数倍に満たない残りは,始点によって値が変わるので省いてはいけない.


%
\section{逆数による置換}\label{節:逆数対称性}

区間の端が $\frac1b$ と $b$ であるときは,\,$x=\frac1t$ でも両端を交換できる.ただし反転 $x=a+b-t$ と異なり,微分は
\[
\dd x=-\frac{1}{t^2}\dd t
\]
になる.関数の $x$ を逆数に変えるだけでは,正しい置換にならない.

特に積分が $\frac{\dd x}{x}$ を含むと,
\[
\frac{\dd x}{x}=-\frac{\dd t}{t}
\]
となり,同じ形に戻る.\,$b>1$ とすると,
\begin{equation}\label{式:逆数置換の対称性}
\int_{\frac1b}^b F(x)\frac{\dd x}{x}
 =\int_{\frac1b}^b F\left(\frac1x\right)\frac{\dd x}{x}
\end{equation}
である.したがって,\,$F(x)+F(\frac1x)=C$ なら
\[
2\int_{\frac1b}^b F(x)\frac{\dd x}{x}
 =C\int_{\frac1b}^b\frac{\dd x}{x}
 =2C\log b
\]
となる.

\begin{example}{逆数による計算}
次の定積分を求めよ.
\[
I=\int_{\frac13}^3\frac{1}{x(1+x^2)}\dd x.
\]
\end{example}

区間の両端が互いに逆数であり,分母に $x$ がある.そこで $F(x)=\frac{1}{1+x^2}$ と見て,\,$\frac{\dd x}{x}$ を残す.
\[
F\left(\frac1x\right)=\cfrac{1}{1+\cfrac{1}{x^2}}
 =\frac{x^2}{1+x^2},
\qquad F(x)+F\left(\frac1x\right)=1.
\]
式\eqref{式:逆数置換の対称性}を使えば
\[
2I=\int_{\frac13}^3\frac{\dd x}{x}
 =\left[\log x\right]_{\frac13}^3=2\log3,
\]
よって $I=\log3$ である.

この例を直接置換しても,
\[
I=\int_3^{\frac13}
\cfrac{1}{\cfrac1t(1+\cfrac1{t^2})}
\left(-\frac{1}{t^2}\right)\dd t
=\int_{\frac13}^3\frac{t}{1+t^2}\dd t
\]
となる.元の被積分関数との和は $\frac1x$ であり,上の計算と一致する.微分因子を忘れると,この和は現れない.

逆数の対称性を,いつもの左右の対称性として見直すこともできる.\,$x=\e^u$ と置くと,\,$\frac{\dd x}{x}=\dd u$,区間は $[-\log b,\log b]$ になる.さらに $x$ を $\frac1x$ に変える操作は,\,$u$ を $-u$ に変える操作になる.数直線上の普通の長さでは $\frac1b$ と $b$ は $1$ から等距離ではないが,対数で測った座標では左右対称なのである.

一般に,端が正で $ab=k$ なら,\,$x=\frac{k}{t}$ は $[a,b]$ の両端を交換する.この場合にも $\frac{\dd x}{x}=-\frac{\dd t}{t}$ が成り立つ.端を見て置換を考えるときは,和が一定か,積が一定かが手がかりになる.


%
\section{和が一になる関数}\label{節:相補重み}

正弦と余弦の節では,反転した積分と加えると分子の和が分母に一致した.原点について対称な区間でも,同じことが起こる分数がある.\,$a>0$,\,$g$ を偶関数,\,$h$ を奇関数とすると,それは
\[
I=\int_{-a}^a\frac{g(x)}{1+\e^{h(x)}}\dd x
\]
の形をしている.指数関数を含む分母を取り除き,\,$\int_0^a g(x)\dd x$ に帰着させよう.

ここでも先に結果を述べ,その後で理由を確かめる.
\begin{equation}\label{式:偶関数と相補重み}
\int_{-a}^a\frac{g(x)}{1+\e^{h(x)}}\dd x
 =\int_0^a g(x)\dd x
\end{equation}
が成り立つ.指数関数の因子を一度も積分せずに,式から取り除けるのだ.

これを示そう.\,$w(x)=\frac{1}{1+\e^{h(x)}}$ と置く.\,$h$ は奇関数なので $h(-x)=-h(x)$ であり,
\[
\begin{aligned}
w(-x)&=\frac{1}{1+\e^{-h(x)}}\\
 &=\frac{\e^{h(x)}}{\e^{h(x)}+1}
\end{aligned}
\]
となる.分母は $w(x)$ と同じで,分子の和 $1+\e^{h(x)}$ が分母に一致するから,
\[
w(x)+w(-x)=1
\]
である.一方,\,$x=-t$ と置くと,\,$g$ が偶関数であることから
\[
I=\int_{-a}^a g(x)w(-x)\dd x
\]
となる.2式を加えると
\[
\begin{aligned}
2I&=\int_{-a}^a g(x)\{w(x)+w(-x)\}\dd x\\
 &=\int_{-a}^a g(x)\dd x\\
 &=2\int_0^a g(x)\dd x
\end{aligned}
\]
を得る.両辺を $2$ で割れば式\eqref{式:偶関数と相補重み}である.

\begin{bookpoint}
\textbf{POINT}\par
$x$ を $-x$ に変えた積分を作ってもとの積分と加え,\,$\frac{1}{1+\e^{h(x)}}+\frac{1}{1+\e^{-h(x)}}=1$ を使って指数関数の因子を消す.
\end{bookpoint}

\begin{example}{指数関数と対称性}
次の定積分を求めよ.
\[
I=\int_{-2}^2\frac{x^4+\cos x}{1+\e^{x^3+x}}\dd x.
\]
\end{example}

\paragraphbf{方針}
分子 $x^4+\cos x$ は偶関数,指数 $x^3+x$ は奇関数である.本節の式で $g(x)=x^4+\cos x$,\,$h(x)=x^3+x$,\,$a=2$ とした場合にあたる.

\paragraphbf{解答}
$I=\frac{32}{5}+\sin2$

\paragraphbf{解法}
$g(x)=x^4+\cos x$ は偶関数で,\,$h(x)=x^3+x$ は奇関数である.\,$x$ を $-x$ に変えると分子はそのまま,指数だけが負になる.したがって
\[
\frac{1}{1+\e^{h(x)}}+\frac{1}{1+\e^{-h(x)}}=1
\]
より
\[
2I=\int_{-2}^2(x^4+\cos x)\dd x
\]
となる.よって
\[
\begin{aligned}
I&=\int_0^2(x^4+\cos x)\dd x\\
 &=\left[\frac{x^5}{5}+\sin x\right]_0^2\\
 &=\frac{32}{5}+\sin2
\end{aligned}
\]
である.分子は区間全体で正である必要はない.式\eqref{式:偶関数と相補重み}に使われているのは,分子が偶関数であることだけである.

\paragraphbf{和が一になる関数の発想}
式\eqref{式:偶関数と相補重み}の出発点は,次の2つの式の比較である.
\[
\frac{1}{1+\e^x},\qquad
\frac{1}{1+\e^{-x}}=\frac{\e^x}{1+\e^x}
\]
2つの分母は等しく,分子の和 $1+\e^x$ が分母に一致するので,2つを足すと $1$ になる.\,$w(x)=\frac{1}{1+\e^x}$ は,対称な2点で値の和が $1$ になる因子なのである.\,$x$ を $-x$ に変えて元の式と加えると,この因子が掛かった部分は $w(x)+w(-x)=1$ に置き換わり,消えてしまう.必要なのは指数が $x$ であることではなく,反転すると指数の符号が変わることである.そのため,指数を奇関数 $h(x)$ に広げても同じ議論が使えた.

\paragraphbf{分子が奇関数の場合}
一方,分子が奇関数のときは,同じ結論にはならない.たとえば $a>0$ として
\[
K=\int_{-a}^a\frac{x}{1+\e^x}\dd x
\]
を式\eqref{式:左右の積分}で左右の組にすると,
\[
\begin{aligned}
K&=\int_0^a x\left(\frac{1}{1+\e^x}-\frac{1}{1+\e^{-x}}\right)\dd x\\
 &=\int_0^a\frac{x(1-\e^x)}{1+\e^x}\dd x
\end{aligned}
\]
となる.被積分関数は $0<x\le a$ で負なので,\,$K<0$ である.「分子の奇関数を対称区間で積分するから $0$」と考えると,分母が偶関数でないことを見落とす.積全体の対称性か,反転後の式との関係を確かめる必要がある.

\paragraphbf{中心を移す場合}
この考えは原点以外にも移せる.区間 $[a,b]$ の中央を $c=\frac{a+b}{2}$ とし,\,$g(a+b-x)=g(x)$,\,$w(x)+w(a+b-x)=1$ なら
\[
\int_a^b g(x)w(x)\dd x=\frac12\int_a^b g(x)\dd x
\]
である.たとえば $w(x)=\frac{1}{1+\e^{x-c}}$ はこの条件を満たす.\,$a+b-x-c=-(x-c)$ だからである.変数を $u=x-c$ に移せば,原点に関する場合と全く同じ証明になる.
%
\section{対称性の使い方}\label{節:定積分の手法選択}

ここまでの方法は,どれも新しい原始関数を増やすものではない.定積分を,値の等しい別の積分に書き換え,2つを組にする方法である.では,与えられた定積分にどの置換を試せばよいのだろうか.最初に区間を見て,次にその区間を保つ変換を考えると,候補を絞れる.

\begin{center}
\begin{tabularx}{0.96\linewidth}{@{}lXX@{}}
\toprule
区間や式の特徴 & 試す変換 & 比べる部分\\
\midrule
$[-a,a]$ & $x\mapsto -x$ & 偶関数と奇関数の部分\\
$[a,b]$ & $x\mapsto a+b-x$ & 和が定数になるか,一次因子が消えるか\\
$\left[0,\frac{\pi}{2}\right]$ & $x\mapsto\frac{\pi}{2}-x$ & 正弦と余弦が交換されるか\\
長い区間と周期関数 & 整数周期だけ平行移動 & 完全な周期と残りの区間\\
$\left[\frac1b,b\right]$ と $\frac{\dd x}{x}$ & $x\mapsto\frac1x$ & 逆数にした関数との和\\
\bottomrule
\end{tabularx}
\end{center}

候補が見つかっても,対称性だけで最後まで解けるとは限らない.その場合は,対称性によって厄介な因子を1つ除き,残った部分を置換積分や部分積分に帰着させればよい.

\begin{bookpoint}
\textbf{POINT}\par
区間を保つ変換で $x$ などの厄介な因子を先に除き,残った積分を置換積分・部分積分・部分分数分解で計算する.
\end{bookpoint}

\begin{example}{反転と置換}
次の定積分を求めよ.
\[
I=\int_0^\pi\frac{x\sin x}{1+\sin^2x}\dd x.
\]
\end{example}

\paragraphbf{方針}
$\frac{\sin x}{1+\sin^2x}$ は $x$ を $\pi-x$ に変えても同じなので,反転で $x$ を $\frac{\pi}{2}$ に置き換える.残りは $u=\cos x$ と置き,部分分数分解する.

\paragraphbf{解答}
$I=\frac{\pi}{\sqrt2}\log(\sqrt2+1)$

\paragraphbf{解法}
$g(x)=\frac{\sin x}{1+\sin^2x}$ は $g(\pi-x)=g(x)$ を満たす.式\eqref{式:対称な重みの一次因子}より
\[
I=\frac{\pi}{2}J,\qquad
J=\int_0^\pi\frac{\sin x}{1+\sin^2x}\dd x
\]
である.\,$J$ では分子の $\sin x\dd x$ が $-\dd(\cos x)$ に対応する.\,$u=\cos x$ と置くと,\,$x=0,\pi$ に対して $u=1,-1$ であり,\,$1+\sin^2x=2-u^2$ だから
\[
J=\int_{-1}^1\frac{1}{2-u^2}\dd u
\]
となる.この分母は $(\sqrt2-u)(\sqrt2+u)$ に分かれるので,
\[
\frac{1}{2-u^2}
 =\frac{1}{2\sqrt2}\left(\frac{1}{\sqrt2-u}+\frac{1}{\sqrt2+u}\right)
\]
である.\,$-1\le u\le1$ では $\sqrt2-u$,\,$\sqrt2+u$ がともに正なので,
\[
\begin{aligned}
J&=\left[\frac{1}{2\sqrt2}
\log\frac{\sqrt2+u}{\sqrt2-u}\right]_{-1}^1\\
 &=\frac{1}{2\sqrt2}\left(\log\frac{\sqrt2+1}{\sqrt2-1}-\log\frac{\sqrt2-1}{\sqrt2+1}\right)\\
 &=\frac{1}{\sqrt2}\log\frac{\sqrt2+1}{\sqrt2-1}\\
 &=\frac{1}{\sqrt2}\log(\sqrt2+1)^2\\
 &=\sqrt2\log(\sqrt2+1)
\end{aligned}
\]
となる.4行目では $(\sqrt2+1)(\sqrt2-1)=1$ より $\frac{\sqrt2+1}{\sqrt2-1}=(\sqrt2+1)^2$ を使った.したがって
\[
\begin{aligned}
I&=\frac{\pi}{2}\cdot\sqrt2\log(\sqrt2+1)\\
 &=\frac{\pi}{\sqrt2}\log(\sqrt2+1)
\end{aligned}
\]
を得る.

\paragraphbf{部分積分から始めた場合}
もし最初に部分積分を使うなら,\,$\frac{\sin x}{1+\sin^2x}$ の原始関数が必要になり,さらにその対数を積分する項が現れる.反転によって $x$ を除いたため,必要な計算が1度の置換と部分分数分解で済んだのである.同じ式でも,不定積分では「微分の逆を作れるか」を見る.定積分ではそれに加えて,「区間を保つ置換で2つの式を補い合わせられるか」を見るとよい.
%
\section*{章末問題}
\begin{enumerate}[label=\textbf{\thechapter.\arabic*},leftmargin=*]
\item\label{問:奇偶の混合}
次の定積分を求めよ.
\[
\int_{-2}^2\left(\frac{x^3}{1+x^2}+2x^2-\cos x\right)\dd x.
\]
\item\label{問:一次重み反転}
反転を利用して,\,$\int_0^\pi x\sin^2x\dd x$ を求めよ.
\item\label{問:三角相補}
$\int_0^{\frac{\pi}{2}}\frac{\cos^6x}{\sin^6x+\cos^6x}\dd x$ を求めよ.分母が $0$ にならないことも確かめよ.
\item\label{問:周期絶対値}
$\int_{\frac{\pi}{3}}^{\frac{10\pi}{3}}|\cos x|\dd x$ と
$\int_{\frac{\pi}{3}}^{\frac{10\pi}{3}}\cos x\dd x$ をそれぞれ求め,値が異なる理由を述べよ.
\item\label{問:逆数四乗}
$\int_{\frac14}^{4}\frac{1}{x(1+x^4)}\dd x$ を,逆数の置換を用いて求めよ.
\item\label{問:指数相補}
$\int_{-1}^1\frac{3x^2+1}{1+\e^{x^5}}\dd x$ を求めよ.
\item\label{問:中心を移す相補}
$\int_2^6\frac{(x-4)^2+2}{1+\e^{x-4}}\dd x$ を求めよ.
\item\label{問:二次重みの限界}
$I=\int_0^\pi x^2\sin x\dd x$ とする.反転によって
$I=\int_0^\pi(\pi-x)^2\sin x\dd x$ と書けるが,この二式を加えるだけで $I$ は決まらないことを説明せよ.その後,部分積分を用いて $I$ を求めよ.
\end{enumerate}

\smallskip\noindent 解答は\pageref{解答:第8章}ページ.

%
